% Generated by code/interviews-1/03-speed-drills/python/gen_drill.py. Do not edit.
\ivfamily{Fractions as percentages}{2, 3, 13, 20, 21, 24, 39, 40, 41, 47}{Know the eighths and the sixteenths. For a denominator of 20, 25 or 50, scale it up to 100: $\tfrac{7}{25} = \tfrac{28}{100}$. For 40, take half of the value over 20.}
\ivfamily{Decimals as fractions}{26, 27, 28, 30, 33, 36, 37, 45}{Read the decimal over a power of 10, then cancel: $0.375$ is $\tfrac{375}{1000}$, which is $\tfrac38$. Most of them are multiples of $\tfrac18$, $\tfrac1{20}$, $\tfrac1{25}$ or $\tfrac1{40}$.}
\ivfamily{Percentage change}{1, 4, 6, 8, 29, 32, 34, 43, 51, 53}{Add or subtract the percentage of the base, or multiply at once: $15\%$ up is $\times 1.15$, and $20\%$ down is $\times 0.8$.}
\ivfamily{As a percentage of}{11, 15, 19, 22, 23, 25, 35, 52}{Cancel first, then scale to 100: $36$ of $120$ is $\tfrac{3}{10}$, or $30\%$.}
\ivfamily{Percentages of any number}{5, 16, 17, 18, 31, 38, 46, 49}{Take 1\% or 10\% and build up, or swap the two: $p\%$ of $n$ equals $n\%$ of $p$.}
\ivfamily{Fractions}{7, 9, 10, 12, 14, 42, 44, 48, 50, 54}{Use the smallest common denominator. To divide, multiply by the upside-down fraction, and cancel first.}
